\documentclass[10pt, conference, letterpaper]{IEEEtran}
\IEEEoverridecommandlockouts

\usepackage{cite}
\usepackage{amsmath,amssymb,amsfonts}
\usepackage{graphicx}
\usepackage{booktabs}
\usepackage{multirow}
\usepackage{url}
\usepackage{textcomp}
\usepackage{xcolor}

\newcommand{\obs}[1]{\vspace{2pt}\noindent\textbf{Observation~#1.}}

\begin{document}

\title{SPLINTER: SPLINe TransfER for Auditable Cross-Network Intrusion Detection}

% INFOCOM submissions are double-blind: keep the block below anonymized for
% submission and restore author names for the camera-ready version.
\author{\IEEEauthorblockN{Anonymous Author(s)}
\IEEEauthorblockA{Submission to IEEE INFOCOM 2026}}

\maketitle

\begin{abstract}
Network intrusion detection systems (NIDS) trained on one network do not transfer to another, because flow feature distributions differ across deployments. Two remedies are common: fine-tuning or domain adaptation on target traffic, and a separate model per site. Fine-tuning needs labelled target traffic and overwrites the source weights, so detection at the original site is no longer guaranteed; per-site training multiplies storage and maintenance; and neither explains what the adaptation changed. We propose activation-gated spline adaptation (AGSA), which exploits the locality of Kolmogorov–Arnold networks (KANs): each learned function is a univariate spline attached to a named flow feature, so adaptation is confined to the splines the target traffic activates. AGSA freezes the source weights, fits an overlay per site, and records activation statistics per overlay. Spline-support routing (SSR) merges these overlays into one model that identifies which network it is observing within a few dozen flows, then applies that site's curves. The consolidated model reaches a mean F1 of 0.82 to 0.90 across domains, stores no traffic, needs no further training, and remains bit-identical to the source model on source traffic. AGSA provides built-in explainability by using feature-specific curves that make both its detection decisions and model adaptations auditable.
\end{abstract}

\begin{IEEEkeywords}
network intrusion detection, test-time adaptation, continual learning,
domain shift, Kolmogorov--Arnold networks, NetFlow
\end{IEEEkeywords}

\section{Introduction}\label{sec:intro}
Machine-learning network intrusion detection systems (NIDS) routinely
report F1 scores above $0.95$ on the network they were trained on. In practice, a detector is trained where
labels exist, such as a lab testbed, a vendor's corpus, or a previously
instrumented site, and runs where they do not: a different enterprise
segment, an IoT deployment, a newly onboarded customer network.
Cross-dataset evaluations
\cite{apruzzese2022cross,layeghy2023generalisability} have established
that the resulting gap is large. What prior work does not provide is
the recovery mechanism a network operator needs. Given a trained model,
an analyst willing to label about a hundred flows, and an unlimited
unlabelled traffic stream, the task is to restore detection quality
online, without full retraining, without a labelled target corpus, and
without destroying the model's value on the network it came from or on
the other networks its operator serves.

This paper builds and evaluates that mechanism. The framework adapts a
pre-trained detector to a new network from a 100-flow labelled pool,
about $10^{-4}$ of the target data, together with the live unlabelled
stream. Every stream batch is predicted \emph{before} it is used for
adaptation \cite{wang2021tent}, so reported performance is what a
deployed system would have delivered. The framework rests on three
design decisions, each forced by the security setting rather than
inherited from test-time adaptation (TTA) as practiced in computer
vision. First, adaptation is anchored: a class-weighted supervised
loss on the small pool, together with a prior-matching guard, prevents
the majority-class collapse that entropy-only adaptation
\cite{wang2021tent} suffers on traffic that is between 5\% and 39\%
attacks. Second, adaptation is selective: only the last encoder layer
and head train, which recovers $91\%$ of the gain of full fine-tuning
at a fifth of the parameters while preserving the frozen structure that
retention, consolidation, and auditability build on; the
norm-parameter recipe imported from vision TTA is, by contrast, the
weakest choice we tested. Third, adaptation is confined:
activation-gated spline adaptation (AGSA) uses the locality of
Kolmogorov--Arnold networks (KANs) \cite{liu2024kan}, whose learned
functions are local B-spline expansions. AGSA tracks which spline
basis regions the target traffic activates and masks gradients to the
rest, freezing source knowledge wherever the target provides no
evidence. This form of retention is not expressible in an architecture
whose every weight affects every input.

We evaluate on every ordered pair of three NetFlow benchmarks that share
an identical 49-feature schema \cite{sarhan2022netflow}, across three
architectures (KAN, 1-D CNN, TabTransformer
\cite{huang2020tabtransformer}), in the standard episodic setting and
in a fleet setting in which independently adapted models must be
consolidated. This paper makes the following contributions.

\begin{itemize}
  \item \textbf{A few-shot test-time adaptation framework for
  cross-network NIDS}, to our knowledge the first, combining a weighted
  supervised anchor, an anti-collapse prior, confidence-filtered
  entropy, stream reconstruction, and selective updates, under a
  predict-then-adapt protocol (Section~\ref{sec:approach}).
  \item \textbf{AGSA}, a retention mechanism intrinsic to KANs that
  confines spline updates to target-activated basis regions
  (Section~\ref{sec:method-agsa}), costs no target accuracy where it
  gates, and produces activation statistics that serve as domain
  fingerprints; and \textbf{SSR}, which turns those fingerprints into a
  router that consolidates per-site spline overlays into one fleet
  model with no training and no stored traffic (Section~\ref{sec:ssr}),
  to our knowledge the first consolidation-by-routing result for NIDS.
  \item \textbf{A systematic cross-network study} of 18
  source$\to$target runs, trainable-subset, loss-term, and label-budget
  ablations, and a weight-space consolidation study, summarized in
  eight observations an operator can act on. Zero-shot transfer
  collapses to mean F1 $0.23$, in one case with an inverted decision
  boundary; the framework restores mean F1 to $0.82$ with the KAN, the
  best mean and the best worst case of the three architectures; and
  the consolidated fleet model reaches an all-domain
  mean of $0.82$ to $0.90$, exceeding a 300-label joint model in two of
  three configurations while using no labels at consolidation time.
  \item \textbf{Auditable adaptation.} KAN spline curves before and
  after adaptation show what changed, and identify from the source
  model alone why one benchmark transfers worst
  (Section~\ref{sec:splines}).
\end{itemize}

\section{Related Work}\label{sec:related}
\textbf{Test-time adaptation.} TENT \cite{wang2021tent} adapts
normalization parameters by entropy minimization;
Schneider et al.~\cite{schneider2020bn} adapt batch-norm statistics; TTT
\cite{sun2020ttt} attaches a self-supervised head whose loss can be
optimized on unlabelled test data. Continual TTA methods add drift
control: CoTTA \cite{wang2022cotta} uses stochastic restore toward
source weights, EATA \cite{niu2022eata} filters unreliable samples and
adds an anti-forgetting regularizer, and SAR \cite{niu2023sar}
stabilizes adaptation in wild, temporally correlated streams. All of
these were designed for, and evaluated on, vision corruptions with
balanced classes. We show that the imbalanced security setting changes
the conclusions qualitatively: entropy-only adaptation is unsafe, and
norm-parameter updates are the wrong subset. We therefore build the
framework around a supervised few-shot anchor and reuse the filtering
and restore ideas as components.

\textbf{Cross-network NIDS evaluation and drift.} Apruzzese et
al.~\cite{apruzzese2022cross} and Layeghy and
Portmann~\cite{layeghy2023generalisability} document severe
cross-dataset degradation of ML-NIDS. The standardized NetFlow feature
sets of Sarhan et al.~\cite{sarhan2022netflow} were created to enable
such studies by removing schema mismatch as a confound, and we build
directly on their v3 conversions. INSOMNIA
\cite{andresini2021insomnia} addresses temporal drift \emph{within} one
network through incremental active learning. Our work addresses the
complementary problem of recovery \emph{across} networks, including
consolidating a multi-network fleet, which to our knowledge has not
previously been studied for NIDS.

\textbf{Continual learning.} Catastrophic forgetting
\cite{mccloskey1989catastrophic} is classically mitigated by
weight-space penalties (EWC \cite{kirkpatrick2017ewc}) or experience
replay \cite{rolnick2019replay}. AGSA takes a third, structural route:
it partitions the spline support instead of penalizing weight motion or
storing data. Weight-space model merging, in the form of soups
\cite{wortsman2022soups} and task arithmetic \cite{ilharco2023task},
promises to consolidate separately adapted models without data. We show
that static fusion fails on adapted NIDS for a measurable reason and
that routing on the KAN's own grid fingerprints succeeds
(Section~\ref{sec:merge}).

\textbf{Kolmogorov--Arnold networks.} KANs \cite{liu2024kan} replace
fixed node nonlinearities with learnable edge functions parameterized as
B-spline expansions. Early work emphasizes accuracy and symbolic
interpretability on scientific tasks. We are, to our knowledge, the
first to exploit the \emph{locality} of the spline parameterization for
test-time adaptation and continual retention, and the first systematic
evaluation of KANs under distribution shift in a security setting.

\section{Background: Kolmogorov--Arnold Networks}
\label{sec:kan}
A KAN layer with $n_{\text{in}}$ inputs and $n_{\text{out}}$ outputs
computes
\begin{equation}
y_j=\sum_{i=1}^{n_{\text{in}}}\phi_{ji}(x_i),\qquad
\phi_{ji}(x)=w^{b}_{ji}\,\sigma(x)+\sum_{k=1}^{G}c_{jik}\,B_k(x),
\label{eq:kan}
\end{equation}
where every edge carries its own learnable univariate function: a
global residual path $w^{b}_{ji}\sigma(\cdot)$ plus a spline expansion
in $G$ B-spline basis functions $B_k$ on a fixed grid over $[-1,1]$.
The property our framework relies on is that each coefficient
$c_{jik}$ is local. Because $B_k$ has compact support, $c_{jik}$
influences the function only for inputs $x_i$ that fall in one
sub-interval of that feature's range.

For an intrusion detector that moves between networks, this locality
has three uses, and the framework draws on all of them.
(a)~\emph{Confinable adaptation.} A network exercises each feature
over its own characteristic region, determined by its port mix, its
RTT regime, and its window-size population. Gradients for coefficients
outside that region carry no signal, only noise, which Adam's
per-coordinate normalization \cite{kingma2015adam} amplifies into
full-sized updates that corrupt the learned function exactly where the
new network provides no evidence. AGSA masks those updates. An
architecture with global weights cannot support this: in an MLP, CNN,
or attention block, every weight touches every input.
(b)~\emph{Readable adaptation.} Because each $\phi_{ji}$ is a
univariate curve attached to a named flow feature, the question of
what adaptation changed has a direct visual answer
(Fig.~\ref{fig:splinesunsw}), and the per-feature importance
$\sum_k|c_{jik}|$ gives an audit trail a security-operations analyst
can act on.
(c)~\emph{Drift telemetry at no cost.} The per-basis activation mass
that AGSA tracks is a running estimate of where traffic lives in
feature space, and its change between deployments is itself a
domain-shift signal, obtained without an auxiliary detector.

\section{Proposed Approach}
\label{sec:approach}

\subsection{Setting and Notation}\label{sec:setting}
We call the procedure few-shot continual test-time adaptation (CTTA).
A detector $f_\theta$ is pre-trained on a labelled \emph{source}
network and must operate on a \emph{target} network with a different
traffic distribution. At the target, two resources exist. The first is
a labelled \emph{pool} $\mathcal{P}=\{(x_i,y_i)\}_{i=1}^{100}$ of 100
stratified flows, about $10^{-4}$ of the target data and the realistic
product of a brief analyst triage session. The second is the
unlabelled live \emph{stream}, processed once as a sequence of batches
$B_1,\dots,B_T$ with $|B_t|=512$. The protocol is online: batch $B_t$
is first scored with the current parameters,
$\hat{y}_t=f_{\theta_{t-1}}(B_t)$, and only then used to compute
$\theta_t$ (predict-then-adapt \cite{wang2021tent}), so no reported
number benefits from a transductive peek. In the fleet setting, copies
of the same source model adapt at multiple sites in parallel and are
afterwards consolidated into one deployable model
(Section~\ref{sec:ssr}), which is the regime an NIDS fleet operates
in.

\subsection{Adaptation Objective}\label{sec:objective}
Each adaptation step minimizes
\begin{equation}
\mathcal{L} = \underbrace{\mathrm{CE}_w(\mathcal{P})}_{\text{anchor}}
 + \underbrace{\lambda_e\, \mathcal{H}_{\text{filt}}(B_t)
 + \lambda_m\, \mathrm{KL}(\bar{p}\,\|\,\pi)}_{\text{unsupervised}}
 + \lambda_r\, \mathcal{R}(B_t),
\label{eq:loss}
\end{equation}
with $\lambda_e{=}\lambda_m{=}1$, $\lambda_r{=}0.5$, where
$\mathcal{R}(B_t)=|B_t|^{-1}\sum_{x\in B_t}\|\hat{x}-x\|^2$ is the
mean stream reconstruction error. Each term addresses a specific
property of security traffic.

\emph{Anchor.} $\mathrm{CE}_w$ is class-weighted cross-entropy on a
pool batch,
\begin{equation}
\mathrm{CE}_w(\mathcal{P}) = -\frac{1}{|B_{\mathcal{P}}|}
\sum_{(x,y)\in B_{\mathcal{P}}} w_y \log p_\theta(y \mid x),
\label{eq:ce}
\end{equation}
with inverse-frequency weights
$w_c = \min\{1/(2\hat{\rho}_c),\, 20\}$, where $\hat{\rho}_c$ is the
frequency of class $c$ in the pool. The cap keeps a pool with a
handful of attacks from destabilizing the high adaptation learning
rate \cite{lin2017focal}. This term carries the
recovery: the loss ablation (Table~\ref{tab:loss}) shows that removing
it is catastrophic, while removing everything else costs a few points.

\emph{Anti-collapse guard.} With
$\bar{p}=|B_t|^{-1}\sum_{x\in B_t} p_\theta(\cdot \mid x)$ the
batch-marginal prediction and $\pi$ the pool's class prior, the term
$\mathrm{KL}(\bar{p}\,\|\,\pi)$ pins the predicted attack rate to the
prior. On traffic that is between 5\% and 39\% attacks, unsupervised
objectives otherwise drift toward predicting the majority class
everywhere; entropy-only adaptation reaches F1 as low as $0.001$
(Section~\ref{sec:ablate}), confidently benign on everything, which is
the one failure a security system cannot tolerate.

\emph{Filtered entropy.} Let
$H(x)=-\sum_y p_\theta(y \mid x)\log p_\theta(y \mid x)$ and
$E_0 = 0.4\ln 2$. Following EATA \cite{niu2022eata}, only reliable
stream samples $S_t=\{x \in B_t : H(x)<E_0\}$ contribute, weighted by
their confidence:
\begin{equation}
\mathcal{H}_{\text{filt}}(B_t) = \frac{1}{|S_t|} \sum_{x \in S_t}
e^{\,E_0 - H(x)}\, H(x).
\label{eq:ent}
\end{equation}
This extracts signal from the unlabelled stream without letting
unreliable samples near the decision boundary steer the update.

\emph{Stream reconstruction.} All models carry a decoder head trained
with a reconstruction loss at the source; at the target the same loss
on stream batches provides self-supervision in the style of TTT
\cite{sun2020ttt} exactly where the shift lives.

Optimization uses Adam with learning rate $10^{-2}$ and stochastic
restore in the style of CoTTA \cite{wang2022cotta}. After each step,
\begin{equation}
\theta \leftarrow m \odot \theta_{\text{src}} + (1-m)\odot\theta,
\qquad m \sim \mathrm{Bernoulli}(0.01),
\label{eq:restore}
\end{equation}
applied coordinate-wise over the trainable subset, which bounds
cumulative drift over the million-flow stream.

\subsection{Selective Spline Adaptation}\label{sec:selective}
Only a small subset of parameters trains: the norm layers, classifier
head, and last encoder layer, between $12\%$ and $21\%$ of the
parameters, or, in the most restricted variant, the last KAN layer's
spline coefficients alone. Selectivity is not only an efficiency
choice. The first encoder layer's splines are the model's learned
per-feature transforms, the transferable knowledge the source
contributes, and freezing them both preserves that knowledge and keeps
the adaptation readable (Section~\ref{sec:splines}). The subset ablation (Section~\ref{sec:ablate}) confirms the design:
restricting updates to the last KAN layer costs almost nothing, while norm parameters, the canonical choice in vision TTA, are the weakest addition, because tabular flow features do not shift primarily in their first two moments.

\subsection{AGSA: Activation-Gated Spline Adaptation}
\label{sec:method-agsa}
AGSA is the mechanism that locality makes possible. For every spline
coefficient $c_{jik}$ of each trainable KAN layer, AGSA maintains an
exponential moving average of the activation mass observed on target
batches,
\begin{equation}
A_{ik} \leftarrow \gamma\, A_{ik} + (1-\gamma)\,
\frac{1}{|B_t|}\sum_{x \in B_t} \big|B_k(x_i)\big|,
\qquad \gamma = 0.99,
\label{eq:ema}
\end{equation}
and masks the gradient of every coefficient whose mass is below $5\%$
of the maximum on its input dimension:
\begin{equation}
\tilde{\nabla} c_{jik} = \nabla c_{jik}\cdot
\mathbf{1}\!\left[\,A_{ik} \geq 0.05\,
\max_{k'} A_{ik'}\right].
\label{eq:agsa}
\end{equation}
Adaptation is thereby confined to spline regions the target
exercises, which in our runs cover $36\%$ to $47\%$ of coefficients,
while the source function elsewhere is structurally frozen: no penalty
term, no extra loss, and no effect on visited-region fitting. The same
averages $A_{ik}$, normalized per input dimension, are a compact
record of where the site's traffic lives on the grid, the statistic
Section~\ref{sec:ssr} turns into a router.

\subsection{Fleet Consolidation: Spline-Support Routing (SSR)}
\label{sec:ssr}
An NIDS fleet adapts in parallel. Each site $D$ produces an
\emph{overlay}, the last layer's adapted spline tensors ($18{,}432$
parameters, ${\sim}74$\,KB), and, as a by-product of AGSA, an
activation-mass \emph{fingerprint} of its traffic. The overlay
suffices: under the full framework these tensors alone land within
$1.4$ points of the full selective subset (Table~\ref{tab:subset})
while retaining more of the source function (mean post-adaptation
source F1 $0.46$ vs.\ $0.35$), so the $74$\,KB overlay carries
essentially the entire adaptation. Normalizing the activation mass per
input dimension gives a distribution over basis bins,
\begin{equation}
p_D[i,k] = \frac{A^{D}_{ik}}{\sum_{k'} A^{D}_{ik'}},
\label{eq:fp}
\end{equation}
a compact, data-free signature of where site $D$'s traffic lives on
the grid. We record it on the frozen first layer, so fingerprints are
directly comparable across sites, and on the adapted layer. SSR
consolidates the fleet into a single model: the shared backbone, one
overlay per site (the source keeps its own), and a router that scores
a window of $W$ flows by fingerprint log-likelihood,
\begin{equation}
\mathrm{route}(x_{1:W})=\arg\max_{D}\;\sum_{w=1}^{W}\sum_{i,k}
\tilde{B}_{ik}(x_w)\,\log p_D[i,k],
\label{eq:ssr}
\end{equation}
where $\tilde{B}(x_w)$ are the normalized B-spline basis activations
that the forward pass computes in any case, so the router costs no
parameters, no training, and no stored flows. (Other architectures can
route too, given added machinery; Section~\ref{sec:discussion}
outlines this extension. The KAN's distinction is that the density model
comes pre-installed: the binning is the model's own grid, the basis
activations are computed by every forward pass, and AGSA already
accumulates the statistics.) One decision is taken per window, not per
flow. A deployed sensor observes a stream from one network segment, so
windows are domain-pure by construction, and aggregation turns a weak
per-flow signal into a near-certain decision. Because the source
retains its own overlay, the consolidated model is bit-identical to
the source model on source traffic.

\subsection{Per-Batch Procedure}
For each stream batch $B_t$: (1)~predict
$\hat{y}_t = f_{\theta_{t-1}}(B_t)$, which is what is scored;
(2)~update the activation averages of Eq.~\eqref{eq:ema} from the same
forward pass and refresh the masks; (3)~draw a pool batch, compute
Eq.~\eqref{eq:loss}, and backpropagate through the trainable subset
with the gated gradients of Eq.~\eqref{eq:agsa}; (4)~apply the Adam
step and the stochastic restore of Eq.~\eqref{eq:restore}. One
optimizer step per stream batch; nothing is revisited.

\section{Experimental Setup}\label{sec:setup}

\subsection{Datasets and Preprocessing}
We use the v3 NetFlow conversions \cite{sarhan2022netflow} of three
standard benchmarks (Table~\ref{tab:data}): NF-CICIDS2018-v3
\cite{sharafaldin2018cicids} (enterprise traffic), NF-ToN-IoT-v3
\cite{moustafa2021toniot} (IoT/IIoT telemetry-heavy traffic), and
NF-UNSW-NB15-v3 \cite{moustafa2015unsw} (synthetic mixed traffic). All
three expose the \emph{same} 49 usable flow features after dropping
identifiers (addresses, ports, timestamps), which isolates distribution
shift from schema mismatch. From each dataset we draw a stratified
sample of $10^6$ flows; detection is binary and F1 is always computed on
the attack class. Duration and bytes-per-packet are derived features,
clipped at zero before log transforms; heavy-tailed volume and
inter-arrival features are $\log(1+x)$ transformed; all features are
mapped to $[-1,1]$ by a rank-based quantile transform fit on the
source only and reused unchanged on the target. Rank scaling is a
deployment choice. Out-of-range target values saturate at the grid
boundary, and the $[-1,1]$ range matches the KAN spline support, so the
position of a flow on the spline grid coincides with its position in
the data distribution.

\begin{table}[t]
\caption{Datasets (v3 NetFlow conversions, $10^6$-flow stratified samples)
and in-domain test F1 after source pre-training.}
\label{tab:data}
\centering
\begin{tabular}{lccccc}
\toprule
 & & Attack & \multicolumn{3}{c}{In-domain F1} \\
\cmidrule(lr){4-6}
Dataset & Traffic & rate & KAN & CNN & TabTr. \\
\midrule
NF-CICIDS2018-v3 & enterprise & 12.98\% & .968 & .966 & .968 \\
NF-ToN-IoT-v3 & IoT/IIoT & 38.97\% & .929 & .924 & .927 \\
NF-UNSW-NB15-v3 & synthetic mix & 5.44\% & .999 & .998 & .998 \\
\bottomrule
\end{tabular}
\end{table}

\subsection{Architectures and Training}
All models share an encoder--classifier--decoder layout returning
(logits, reconstruction, latent). The \textbf{KAN}
\cite{liu2024kan,efficientkan} encoder is $49{\to}64{\to}32$ with grid
range $[-1,1]$ and $G=8$ basis functions per edge (103{,}810
parameters, 20{,}610 in the selective subset);
the \textbf{CNN} is a compact 1-D convolutional encoder (10{,}963 /
2{,}338); the \textbf{TabTransformer} \cite{huang2020tabtransformer}
tokenizes each feature and applies 4-head self-attention (6{,}147 /
738). The CNN and TabTransformer serve as controls with global weights:
identical pipeline and objectives, no locality. Each model is
pre-trained on its source for 20 epochs (Adam, learning rate $10^{-3}$,
weighted CE plus $0.5\times$ reconstruction); in-domain test F1 ranges
from $0.92$ to $0.999$ (Table~\ref{tab:data}), so nothing that follows
is an artifact of a weak source model. At the target, the stream of
999{,}900 flows yields 1{,}953 online steps.

\subsection{Baselines and Controls}
\emph{Zero-shot} applies the source model unchanged. \emph{Pool-only
fine-tuning} (``few-shot'') trains the same selective subset on the
pool alone with the same optimizer, that is, the framework minus the
stream. Ablations vary the trainable subset and delete loss terms one
at a time under identical budgets. Training \emph{all} parameters on
the same protocol serves as an upper reference. It adds a few points
of target F1 (KAN mean $0.872$ vs.\ $0.818$) but rewrites the frozen
first-layer transforms that retention, fleet consolidation, and the
audit trail of Section~\ref{sec:splines} are built on; the selective
subset recovers $91\%$ of that gain with $20\%$ of the parameters, so
we treat full tuning as a bound rather than a deployable
configuration. All experiments
use a single seed re-applied at the start of every configuration;
paired comparisons, for example AGSA on versus off, share weights and
RNG state so the mechanism under test is the only difference. A complete end-to-end
re-execution of the KAN experiments reproduced every headline aggregate
within $\pm 0.03$ F1 and every qualitative conclusion; the claims in
this paper are those that held across both executions. Fisher-information (EWC
\cite{kirkpatrick2017ewc}) retention was implemented but disabled after
the measured penalty proved negligible ($\leq 8\times10^{-4}$;
Section~\ref{sec:agsa}).

\section{Cross-Network Results}\label{sec:results}

\begin{figure}[t]
\centering
\includegraphics[width=\columnwidth]{fig_trajectory.png}
\caption{Online (predict-then-adapt) F1 over the stream for the KAN
pre-trained on CICIDS2018. Adaptation to ToN-IoT saturates within
${\sim}150$ batches (${\sim}75$k flows); UNSW-NB15, whose zero-shot F1 is
$0.012$, needs ${\sim}750$ batches (${\sim}380$k flows) to plateau.}
\label{fig:traj}
\end{figure}

\input{tab_main}

\obs{1} \emph{A shared feature schema does not transfer, and the
failure can be inverted.} Despite identical features and near-ceiling
source performance, zero-shot F1 averages $0.232$ over the 18 pairs
(Table~\ref{tab:main}), with UNSW-NB15 as target the worst (F1 between
$0.011$ and $0.088$ from the CICIDS2018 and ToN-IoT sources). The
failure is not mere under-confidence. Fig.~\ref{fig:logits} (left)
shows that the CICIDS2018 model assigns higher attack logits to benign
UNSW flows than to actual attacks; the decision boundary is inverted
in the shifted representation. The operational consequence is that a
drift monitor relying on the source model's own confidence would not
merely miss the degradation but would be systematically misled, and
recalibrating the threshold cannot fix a boundary that points the
wrong way.

\obs{2} \emph{One hundred labelled flows plus the stream recover most
of the loss, and the stream is what makes the recovery robust.} The
framework lifts mean F1 from $0.232$ to $0.819$ (KAN), $0.774$ (CNN),
and $0.760$ (TabTransformer), recovering, against the in-domain
ceiling of $0.92$ to $1.0$, the bulk of what the network move
destroyed, with per-pair gains of up to $+0.85$ (Table~\ref{tab:main},
$\Delta$ column). Post-adaptation accuracy averages $0.93$ for the KAN
against $0.91$ and $0.90$ for the controls. For the architectures with
global weights the stream is necessary: it adds $+0.07$ (CNN) and
$+0.25$ (TabTransformer) over pool-only fine-tuning, whose floor
collapses to $0.50$ and $0.32$ on their worst pairs. The KAN is the
exception, and the reason is locality. Each labelled flow adjusts only
the spline regions it lands in, so a 100-flow stratified pool alone
already reaches mean $0.823$, matching CTTA's mean within noise. Even
there the stream matters on the quantities an operator provisions
for: it raises the worst-case pair from $0.729$ to $0.771$,
and Fig.~\ref{fig:traj} shows the inverted-boundary shift
(CIC$\to$UNSW, zero-shot $0.012$) climbing for about $380$k flows, an
improvement no offline pool fit can deliver. The label budget
completes the argument (Table~\ref{tab:pool}). Halving the pool to
$50$ flows costs pool-only fine-tuning roughly twice what it costs
CTTA (KAN means $0.82\!\to\!0.75$ vs.\ $0.82\!\to\!0.78$) and reverses
their order; with 50 labels the stream is worth $+2.7$ points of mean
F1 and up to $+21$ on the severest shift (ToN$\to$UNSW: $0.98$ vs.\
$0.77$), while the KAN's lead over the CNN and TabTransformer persists
($0.78$/$0.74$/$0.66$). The scarcer the labels, the more the stream
matters. At the 100-flow budget the KAN posts both the highest mean
and the highest worst-case pair ($0.771$, vs.\ $0.697$ for the CNN and
$0.591$ for the TabTransformer), and in security deployment the worst
case is the operative number, because a detector is provisioned for
its floor.

\begin{figure*}[t]
\centering
\includegraphics[width=0.82\textwidth]{fig_logits.png}
\caption{Logit margin (attack $-$ benign) distributions on UNSW-NB15 for
the KAN pre-trained on CICIDS2018, before (left) and after (right)
adaptation. Before adaptation the classes are interleaved and partially
inverted, in that attack flows receive lower attack margins than benign
ones. Adaptation changes every flow's margin, so the distributions
shift as well as separate: after adaptation each class migrates to its
correct side of the boundary with larger magnitude, and the horizontal
ranges of the two panels differ accordingly.}
\label{fig:logits}
\end{figure*}

\obs{3} \emph{Adaptation must be paired with retention.} An adapted
model is a specialist. After CTTA, F1 on the \emph{source} test set
falls to $0.35$ on average for the KAN, and the loss is strongly
directional: models pre-trained on UNSW-NB15 forget it almost
completely, while adaptation toward UNSW-NB15 from ToN-IoT retains
$0.71$. The asymmetry mirrors dataset breadth. UNSW-NB15's decision
function occupies a narrow, brittle region that any adaptation
overwrites, whereas broader source distributions such as ToN-IoT share
support with their targets; Section~\ref{sec:splines} traces this
brittleness to a visible property of what the UNSW model learned. A
deployable system therefore needs a retention mechanism alongside
adaptation, here AGSA plus overlay-based consolidation, and
Section~\ref{sec:merge} shows that the consolidated model recovers the
source exactly.

\section{Ablations: Where and What to Adapt}\label{sec:ablate}

\input{tab_subset}
\input{tab_loss}
\input{tab_pool}

\obs{4} \emph{Where you adapt matters more than how much, and locality
is what makes a small update surface sufficient.}
Table~\ref{tab:subset} compares, under the full framework, the
standard subset (norm layers, classifier, and last encoder layer) with
the last encoder layer alone, for all three architectures. Restricting
the update costs the KAN $1.4$ points ($0.818$ to $0.804$), the CNN
$2.0$, and the TabTransformer $11.1$; the KAN's last layer alone
outperforms every configuration of the other architectures. A
pool-only sweep over finer subsets points the same way: the last KAN
layer is the best single choice (mean $0.830$), and norm parameters,
the canonical choice of TENT-style vision TTA
\cite{wang2021tent,schneider2020bn}, are the weakest addition tested,
with splines+norm ($0.761$) falling below spline-only ($0.772$). The
explanation is structural. Tabular flow features, unlike corrupted
images, do not shift primarily in their first two moments, so
recalibrating affine normalization fixes little and adds interference;
the useful degrees of freedom are the function-shape parameters
nearest the head, and only the KAN concentrates them there. This is
what makes the per-site overlays of Section~\ref{sec:ssr} possible.


\obs{5} \emph{The labelled anchor does the work; entropy minimization
alone is catastrophic under imbalance.} Table~\ref{tab:loss} deletes
terms from Eq.~\eqref{eq:loss} under the spline-only subset. Pure
confidence-filtered entropy minimization, the TENT-like variant,
collapses on five of six pairs and bottoms out at F1 $0.001$ on
CIC$\to$UNSW. With attack rates between 5\% and 39\%, the entropy
gradient is dominated by the majority class and the model converges to
confidently predicting benign everywhere, which is the failure a
security system cannot tolerate and one that balanced-class vision
benchmarks never surface. Any variant containing the 100-flow anchor
lands between $0.72$ and $0.76$ mean F1. The full objective is best
($0.765$ mean, and $+0.10$ over CE-only on the severe CIC$\to$UNSW
shift), with the unsupervised terms adding $+0.04$ mean over the anchor
alone, a consistent gain but an order of magnitude smaller than the
anchor's contribution. For NIDS, test-time adaptation in the strictly
unsupervised sense \cite{wang2021tent} is not merely suboptimal; it is
unsafe. This is the empirical footing for the framework's anchor-first
design.

\subsection{AGSA in the Episodic Setting}\label{sec:agsa}
\obs{6} \emph{Structural gating is free, and its statistics are the
fleet's telemetry.} We compare spline-only adaptation with the AGSA
gate on versus off under identical weights and RNG on all six KAN
pairs. The target stream activates only $36\%$ to $47\%$ of spline
coefficients, so the gate freezes more than half of the spline
support, yet mean target F1 is fully preserved ($0.804$ gated vs.\
$0.799$ ungated), and on the two severest shifts the gate gains
$+3.0$ and $+3.7$ points (CIC$\to$UNSW and ToN$\to$UNSW). This
confirms the design premise: dormant-region coefficients contribute
nothing to target fitting, so protecting them costs nothing. The
measured EWC penalty after gated adaptation is
$\leq 8\times 10^{-4}$, meaning that in Fisher-weighted terms the
source function barely moves. The gate's by-product is as valuable as
its protection: the per-domain activation masses it accumulates, with
distinct per-domain signatures, are exactly the fingerprints that let
Section~\ref{sec:merge} consolidate a fleet by routing.

\section{Reading Adaptation Off the Splines}\label{sec:splines}

\begin{figure*}[t]
\centering
\includegraphics[width=0.79\textwidth]{fig_splines_unsw.png}
\caption{KAN spline activations before (top two rows) and after (bottom
two rows) adaptation for UNSW-NB15$\to$CICIDS2018: the four most
important input features (first encoder layer) and hidden units (last
encoder layer). The UNSW-pre-trained first layer is dominated by
TTL-derived features (\texttt{MAX\_TTL}, \texttt{MIN\_TTL}) with
near-linear, low-importance curves ($1.43$, an order of magnitude below
the $15.0$ a ToN-IoT-trained model assigns its top feature): the source
task was solved by a low-complexity shortcut. Selective adaptation
leaves the first layer untouched (rows 1 vs.\ 3); the last layer must
compensate, its top importance jumping from $0.23$ to $5.08$ and its
functions becoming strongly non-linear (rows 2 vs.\ 4). The spline view
thus records \emph{what} adaptation changed and \emph{diagnoses},
before any target data arrives, why this source transfers worst
(Observations~1 and~3).}
\label{fig:splinesunsw}
\end{figure*}

The spline view turns the paper's quantitative results into a
mechanism, a capability unique to the locality-based design. For a
healthy source (ToN-IoT), the same analysis, which we summarize in
text for space, shows what selective adaptation should look like: the
first layer's learned feature
transforms, rich high-importance curves on TCP window (importance
$15.0$), application protocol, and flag features, are preserved
verbatim, and adaptation is absorbed entirely by last-layer
recombination of hidden units. The frozen curves are not a side effect
but the point: they are the transferable knowledge the source model
contributes, and the reason recovery from ToN-IoT is cheap.

Fig.~\ref{fig:splinesunsw} (UNSW-NB15 source) shows the pathological
case. The first layer barely learned anything beyond gently sloped
functions of TTL-derived features, a known artifact-prone shortcut in
UNSW-NB15, sufficient for $0.999$ in-domain F1 but nearly useless
elsewhere. Under adaptation the only place capacity exists is the last
layer, which contorts (importance $0.23\to 5.08$, strongly oscillatory
curves) to rebuild discriminative structure from weak inputs. This
single picture explains the cluster of UNSW anomalies in
Table~\ref{tab:main}: the worst zero-shot transfer out of UNSW, and the
near-total forgetting of UNSW after adaptation away from it. The narrow
shortcut is overwritten immediately because little transferable
structure holds it in place. For the CNN and TabTransformer the same
pathology can only be inferred indirectly; the KAN exposes it as an
artifact of the source data that an operator could catch before
deployment. In a domain where analysts must justify why a detector
fired, and why it was retrained, this readability is a deployment
requirement rather than a convenience.

\section{Fleet Consolidation: Merging Fails, Routing Succeeds}
\label{sec:merge}

\input{tab_merge}
\input{tab_ssr}

The episodic grid adapts one model to one network. A fleet operator
faces the next problem: several copies of the source model have
adapted at several sites, and operations wants one model that serves
them all, including the source network, without shipping traffic
between sites. Model merging, in the form of soups
\cite{wortsman2022soups} and task arithmetic \cite{ilharco2023task},
promises exactly this with no data at all, so it is the natural first
candidate. We test it two ways on the KAN with source CICIDS2018.

\emph{Interpolation as a retention dial.} For a single adapted model,
$\theta(\alpha)=\theta_{\text{src}}+\alpha(\theta_{\text{adapted}}
-\theta_{\text{src}})$ sweeps a source/target trade-off using nothing
but the two checkpoints; only the two adapted spline tensors differ.
For the moderate shift (CIC$\leftrightarrow$ToN) this dial is
effective: $\alpha{=}0.5$ scores $0.893$ on the source and $0.692$ on
the target, buying back $50$ points of source F1 for $9$ points of
target, post hoc and without data. Under the inverted shift
(CIC$\leftrightarrow$UNSW) no useful intermediate exists: target F1
stays below $0.08$ until $\alpha\geq 0.75$. Interpolation works where
domains share support; it cannot average across an inverted boundary.

\obs{7} \emph{Static merging fails, and the activation masks explain
why.} Table~\ref{tab:merge} merges two copies of the source model
adapted in parallel to ToN-IoT and UNSW-NB15. Every static fusion
collapses on at least one domain, with all-domain means between $0.45$
and $0.49$. The AGSA activation masks of the two targets overlap
$77\%$ (Jaccard), so the parallel adaptations compete for the same
spline regions, and any weight arithmetic in those contested regions
averages away the more delicate solution. Locality softens the damage;
the AGSA-routed static merge preserves the source best ($0.761$ vs.\
$0.744$ uniform, with dormant regions provably identical to the
source). But no single coefficient tensor can hold two conflicting
functions on the same support.

\obs{8} \emph{Routing at deployment granularity turns failed merging
into lossless consolidation, in every configuration.} If the domains
cannot share one function, the consolidated model can instead choose
the function per window, using the SSR router of Eq.~\eqref{eq:ssr}.
Per-flow routing is genuinely impossible: accuracy stalls between
$69\%$ and $86\%$ across configurations, because single flows,
especially benign ones, are domain-ambiguous, and the per-flow SSR row
of Table~\ref{tab:merge} inherits that ceiling. The signal aggregates
fast, however. In all three source configurations
(Table~\ref{tab:ssr}) the router is at least $99.7\%$ correct at
16-flow windows and exact from $64$, and with one decision per
evaluation batch of at least 512 flows it routed all $2{,}349$ batches
across the three configurations correctly. Consolidation is therefore
lossless: every domain receives exactly its own expert's F1, and on
source traffic the consolidated model is bit-identical to the source
model ($0.968$/$0.929$/$0.999$, matching the in-domain scores to four
decimals). The all-domain mean reaches $0.861$/$0.895$/$0.821$ with no
stored traffic and no consolidation training, and on streams that
switch network mid-flight ($100$k flows per segment) the router tracks
every change within at most one window. A single-model alternative we
also evaluated, which carries one detector through the three networks
sequentially while replaying each site's 100-flow pool with AGSA
enabled, tops out at endpoint all-domain means of $0.78$ to $0.80$
while requiring stored pools and a serial pass; routing dominates it
on every axis. The joint-pool baseline (Joint$_{300}$ in
Tables~\ref{tab:merge} and~\ref{tab:ssr}) sharpens the division of
labor. Despite consuming all $300$ pool labels at consolidation time,
it trails label-free SSR by $7$ to $11$ points of mean in the two
broad-source configurations and degrades the source by $10$ to $15$
points, because one function must compromise across conflicting
domains. What routing cannot do is repair a weak expert. From the
shortcut-ridden UNSW-NB15 source, the UNSW$\to$CIC overlay itself
reaches only $0.623$ (Section~\ref{sec:splines} explains why), the
consolidated floor inherits it, and there the joint model's labels
partially repair what routing can only reproduce ($0.831$ vs.\ $0.821$
mean): routing preserves experts, and only labels can repair one.
Floors are set by adaptation quality, never by consolidation. The
remaining caveat is inherent. Window routing assumes domain-pure
windows, which holds for a sensor on one segment but not for a tap
aggregating several networks, where a single jointly fitted function
(the Joint$_{300}$ row) is the fallback. This completes the retention
picture: masks own the dormant regions, labels own the shared regions,
and routing owns the fleet case, each mechanism reading the same
spline-locality structure a different way.

\section{Why a KAN Wins This Setting}\label{sec:whykan}
The setting plays to the KAN's structure in four measured ways.
\emph{(1)~Recovery.} The KAN posts the highest mean F1 under the
framework ($0.819$ vs.\ $0.774$ and $0.760$), the highest worst-case
pair ($0.771$ vs.\ $0.697$ and $0.591$), and the strongest pool-only
recovery ($0.823$): because spline updates are local, a hundred
labelled flows reshape exactly the regions target traffic occupies,
while the controls need the stream at any label budget.
\emph{(2)~Capacity where selective adaptation needs it.} Restricting
the update to the last spline layer costs the KAN only $1.4$ points,
against $11.1$ for the TabTransformer (Table~\ref{tab:subset}), so the
cheapest update surface carries essentially the whole adaptation,
which makes low-memory, on-path adaptation plausible on network
appliances.
\emph{(3)~Structural retention.} AGSA has no analogue under global
weights: it is free episodically (Observation~6) and its statistics
are the fingerprints that make lossless consolidation possible
(Observation~8), so one mechanism buys both protection and telemetry.
\emph{(4)~Auditability.} Because each learned function is a curve over
a named feature, the spline view turns adaptation into an inspectable
object and diagnosed a source-data shortcut (the TTL features of
UNSW-NB15) that no aggregate metric of the controls reveals
(Section~\ref{sec:splines}). These advantages are largest exactly in
the regime operators face, few-shot selective adaptation with
retention requirements, and the KAN's ${\sim}10^5$ parameters are
negligible for any network appliance. Locality is the one property the
competing architectures cannot bolt on.

\section{Discussion and Limitations}\label{sec:discussion}
\textbf{What generalizes.} Three findings were stable across all
architectures and directions: zero-shot collapse despite schema
parity, the dominance of a small supervised anchor over any
unsupervised signal, and severe asymmetric forgetting. None of them
depends on the encoder's inductive bias, so we expect them to hold for
other flow-based detectors. The subset, gating, and routing results
(Observations~4 and 6--8) are KAN-specific by construction.

\textbf{Porting SSR beyond the KAN.} Nothing in Eq.~\eqref{eq:ssr} is
KAN-specific except the source of the density: any architecture can
route per-site overlays if it is paired with a per-site density
estimate over the input features, for instance per-feature histograms
fitted on site traffic and scored with the same window vote. The
difference is the cost of that estimator. For the KAN it comes
pre-installed: the binning is the model's own spline grid ($49$
features $\times$ $8$ basis functions, a $392$-bin fingerprint of
about $1.5$\,KB), the activations are computed by every forward pass,
and AGSA accumulates the statistics during adaptation, with no added
parameters and no training. The natural in-model statistics of the
controls are weaker. A CNN's batch-norm moments summarize each channel
by mean and variance, precisely the statistics the subset ablation
found least informative for flow-level shift, and the TabTransformer
exposes no per-feature-region structure at all, so both require an
external estimator that must be binned, fitted per site, and shipped
as separate state. Ported overlays would be smaller ($2{,}338$ and
$738$ parameters), but the fleet would inherit weaker experts (means
$0.774$/$0.760$, floors $0.697$/$0.591$, vs.\ $0.819$/$0.771$ for the
KAN). Quantifying such extensions is future work.

\textbf{Operator's recipe.} (1)~Before moving a detector, inspect what
the source model learned; flat, low-importance first-layer splines are
a transfer red flag (Section~\ref{sec:splines}). (2)~Label about 100
stratified flows before trusting any adaptation. (3)~Adapt the last
encoder layer plus head, not the norm layers, with the pool as anchor,
the marginal-prior guard, and stream reconstruction. (4)~Never deploy
unsupervised entropy adaptation on imbalanced security traffic.
(5)~Adapt with AGSA enabled; the gate is free and its fingerprint is
required later for consolidation. (6)~Archive the source checkpoint;
interpolation is a retention dial for moderate shifts, but checkpoint
arithmetic is not a substitute for routing (Section~\ref{sec:merge}).
(7)~Keep each site's overlay and fingerprint; window-SSR consolidates
the fleet into one self-routing model at no data cost, provided each
sensor's window sees one network at a time.

\textbf{Limitations.} Labels are binarized, so per-attack-family
dynamics are invisible. The three benchmarks are research corpora with
known labelling artifacts \cite{apruzzese2022cross}, although
Section~\ref{sec:splines} shows the methodology detecting exactly such
an artifact. The stratified pool assumes the analyst labels a roughly
representative sample. Each target stream is stationary, whereas real
traffic also drifts within a site \cite{andresini2021insomnia}. The
consolidation study covers one architecture, by design, since SSR
requires the spline grid that only the KAN provides.

\section{Conclusion}\label{sec:conclusion}
When a flow-based intrusion detector moves to a new network its
accuracy collapses, occasionally to inverted decisions, even under an
identical feature schema. We presented a few-shot test-time adaptation
framework that recovers this loss from a hundred labelled flows and
the live stream, restoring mean F1 from $0.23$ to $0.82$ with the KAN
achieving the best mean and the best worst case of the architectures
tested. The statistics its gating mechanism accumulates then turn a
fleet of adapted copies into a single self-routing model that gives
every network its own expert's accuracy, a mean of $0.82$ to $0.90$,
while remaining bit-identical to the source model on source traffic
and storing no traffic. The negative results are equally actionable:
unsupervised entropy minimization is unsafe on imbalanced security
streams, norm-only adaptation transplants poorly from vision, and
static weight merging cannot consolidate what routing can. The
underlying argument is architectural. Locality makes adaptation
confinable, readable, and self-monitoring, and on every aggregate
measure of this study it is the property that wins.

\bibliographystyle{IEEEtran}
\bibliography{references}

\end{document}
